\subsection{The steps, the samples, and the conventions}\label{sec:gates}

The record's steps, the outcome names, and the counting conventions are
set out once, here, and every later section uses them.

A US trademark passes through a sequence of dated, high-stakes steps, and
each step leaves a record. An application is filed with a goods/services
description; an examiner reviews it, often issuing office actions the
applicant must answer; if allowed and (for intent-to-use filings) a sworn
statement of use is filed --- the proof that the mark has entered commerce
--- the mark registers. Registration starts a maintenance clock: in years
five to six the owner must file a Section~8 declaration proving continued
use in commerce, with a six-month grace period --- the \emph{five-year
proof of continued use}, the paper's central outcome; around year ten a
combined Section~8 declaration and Section~9 renewal is due, and every ten
years thereafter. A mark that misses a proof is cancelled. Madrid Protocol
registrations (\S66(a) basis) file Section~71 declarations on the same
schedule and are flagged separately throughout. \citet{Graham2013}
describe the administrative data; the event-dated corpus built here adds
the full prosecution history of each case, so a cancellation is assigned
to a specific proof by its date rather than inferred from a status
snapshot.\footnote{Cross-checked against the USPTO status-code dictionary:
700 registered; 701--705 Section~8 accepted; 710 cancelled for Section~8
failure; 800 renewed; 900 expired.}

\begin{table}[h]\centering\small
\caption{The trademark funnel in this corpus. Registration counts are
unique serials; five-year-proof outcomes are event-dated (a Section~8
cancellation at registration age 4.0--8.5 years). Source: the re-parsed
USPTO application backfile, scored as in this section.}
\label{tab:funnel}
\begin{tabular}{lrr}
\toprule
Stage & $n$ & Share \\
\midrule
Debut filings scored, 1995--2018 (owners' first filings) & 3{,}589{,}068 & \\
\quad reached registration & 2{,}066{,}600 & 57.6\% \\
\midrule
All unregistered class-records filed 1995--2018 & 3{,}798{,}318 & \\
\quad never filed statement of use & & 42.6\% \\
\quad stopped responding in prosecution & & 45.3\% \\
\quad removed adversarially & & 4.3\% \\
\midrule
Registrations 2002--2018, scored, unique serials & 3{,}104{,}534 & \\
\quad failed the five-year proof of continued use (event-dated) & & 52.3\% \\
\midrule
Passed the five-year proof, cohorts 2002--2013, status known & 1{,}129{,}724 & \\
\quad failed at the year-ten renewal & & 39.1\% \\
\bottomrule
\end{tabular}
\end{table}

Three named outcomes recur (Table~\ref{tab:funnel}). \emph{Failing to
register}: the application never completed, which happens by the
applicant's own inaction in roughly nine cases of ten
(Section~\ref{sec:val_registration}). \emph{Failing the five-year proof of
continued use}: a registration cancelled for non-use at registration age
4.0 to 8.5 years, which means the product the mark named left commerce
within six years (Section~\ref{sec:val_gate}).
\emph{Reaching the SEC record}: the owner appears in dated SEC events ---
a priced Form~D round, financial reporting, a listing
(Section~\ref{sec:links}).

Outcomes are rates, so every estimate is a shift in a probability and is
stated in percentage points against the base rate it shifts: $+1.4$pp on
a 52.3\% base means 53.7\% against 52.3\%. The basic comparison is the
outcome rate of the fifth of filings with the highest lead minus that of
the fifth with the lowest, with the fifths cut inside a single Nice class
and year so that no difference between industries or between years
enters; regression tables add drafting and filer controls and cluster
errors on the owner. Where a continuous version is given, it is the
change in the outcome for a one-standard-deviation change in the axis, in
percentage points.

Counting differs by table. A filing may claim several Nice classes at
once, so a \emph{class-record} counts it once per class while a
\emph{registration} counts it once by serial number; a \emph{debut
filing} is an owner's first, one per owner. A \emph{cohort} is the set of
registrations issued in a calendar year; where a split is instead by
filing year, the text says so.

Two windows bound the samples, each set by an institutional fact. Scoring
runs 1995--2019: a full five-year past reference does not exist before
1990 and a full forward reference does not exist after 2020, and the 1995
cut is about two years more conservative than the measured burn-in
requires (Section~\ref{sec:burnin}). Five-year-proof cohorts run
2002--2018: cancellations post at about registration age six and a half,
the record ends in April 2026, so 2018 is the last cohort whose window
has essentially elapsed, and 2002 is the conventional lower cut for this
corpus --- a bound that matters, because carried back to the first
scoreable cohorts the lead penalty was larger than it has ever been
since, as the companion paper shows.

% Neighbours/nomenclature moved to Section 2 (rp_related) for the RP build.
